\section{Find Median from a Data Stream}
\label{sec:heap-median-stream}

\problemstatement{Design a structure supporting $\textit{addNum}(x)$,
which adds an integer to a growing stream, and $\textit{findMedian}()$,
which returns the median of all numbers added so far, at any point,
efficiently after every insertion.}

\begin{patternbox}
\emph{``Running median, computed after every new element of a stream''}
is the canonical trigger for the \textbf{two-heap} pattern flagged in
Section~\ref{sec:heap-intro}: a median needs access to both ``the largest
of the lower half'' and ``the smallest of the upper half'' of the data
seen so far, and no single heap can offer both -- so we keep two heaps of
opposite polarity, each covering one half, and keep them balanced.
\end{patternbox}

\subsection*{Understanding the problem carefully}

If we maintain the invariant that all elements in a \textbf{max-heap}
called \textit{lower} are $\le$ all elements in a \textbf{min-heap} called
\textit{upper}, and that the two heaps' sizes differ by at most $1$, then:

\begin{itemize}
  \item If the total count is odd, the median is the top of whichever
        heap has one more element than the other.
  \item If the total count is even, the median is the average of
        $\textit{lower}$'s top (the largest of the lower half) and
        $\textit{upper}$'s top (the smallest of the upper half).
\end{itemize}

The entire problem reduces to maintaining this ``balanced split into a
max-heap lower half and a min-heap upper half'' invariant after every
insertion.

\begin{ideabox}[Insertion Rule]
On $\textit{addNum}(x)$: push $x$ into $\textit{lower}$ (the max-heap).
Then, to preserve the ``everything in \textit{lower} is $\le$ everything in
\textit{upper}'' invariant, pop $\textit{lower}$'s maximum and push it into
$\textit{upper}$ (this guarantees the largest element currently in
\textit{lower} is never larger than anything already confirmed to belong
in the upper half). Finally, to preserve the size-balance invariant: if
$\textit{upper}$ now has more elements than $\textit{lower}$, pop
$\textit{upper}$'s minimum and push it back into $\textit{lower}$.
\end{ideabox}

\subsection*{Why pushing into \textit{lower} first, then shuffling, is correct}

Always inserting $x$ into $\textit{lower}$ first (regardless of whether
$x$ is actually a ``low'' or ``high'' value) and then immediately moving
$\textit{lower}$'s current maximum into $\textit{upper}$ is a clean way to
guarantee correctness without an explicit case split on where $x$ belongs:
after this shuffle, $\textit{upper}$'s new minimum is exactly
$\min(x, \text{old } \textit{upper}\text{'s minimum})$, and one can verify
that whichever of these is smaller is correctly placed at the boundary,
while $\textit{lower}$'s new maximum is whatever was second-largest among
$\{x\} \cup \textit{lower}\text{'s old elements}$, preserving
$\textit{lower} \le \textit{upper}$ throughout. The subsequent
size-rebalancing step (moving one element back if \textit{upper} becomes
too large) maintains the size-difference-at-most-$1$ invariant, which is
exactly what makes the parity-based median formula above always correct.

\subsection*{Algorithm}

\begin{enumerate}
  \item \textsc{AddNum}($x$):
  \begin{enumerate}
    \item Push $x$ onto $\textit{lower}$ (max-heap).
    \item Pop $\textit{lower}$'s max; push it onto $\textit{upper}$
          (min-heap).
    \item If $|\textit{upper}| > |\textit{lower}|$: pop $\textit{upper}$'s
          min; push it back onto $\textit{lower}$.
  \end{enumerate}
  \item \textsc{FindMedian}():
  \begin{enumerate}
    \item If $|\textit{lower}| > |\textit{upper}|$: return
          $\textit{lower}$'s top.
    \item Else: return $(\textit{lower}\text{'s top} +
          \textit{upper}\text{'s top}) / 2$.
  \end{enumerate}
\end{enumerate}

\begin{algorithm}[H]
\caption{Median Maintainer (Two Heaps)}
\begin{algorithmic}[1]
\Procedure{AddNum}{$x$}
  \State push $x$ onto \textit{lower} (max-heap)
  \State push (pop \textit{lower}) onto \textit{upper} (min-heap)
  \If{$|\textit{upper}| > |\textit{lower}|$}
    \State push (pop \textit{upper}) onto \textit{lower}
  \EndIf
\EndProcedure
\Procedure{FindMedian}{}
  \If{$|\textit{lower}| > |\textit{upper}|$} \Return top of \textit{lower}
  \Else \Return $(\text{top of } \textit{lower} + \text{top of } \textit{upper}) / 2$ \EndIf
\EndProcedure
\end{algorithmic}
\end{algorithm}

\subsection*{Dry run}

\dryrun{Add $5, 15, 1, 3$ in order, checking the median after each.}

\begin{itemize}
  \item Add $5$: push to \textit{lower} $=\{5\}$. Pop max ($5$), push to
        \textit{upper} $=\{5\}$. Now $|\textit{upper}|=1 >
        |\textit{lower}|=0$: pop \textit{upper}'s min ($5$), push to
        \textit{lower}. \textit{lower}$=\{5\}$, \textit{upper}$=\{\}$.
        Median: $|\textit{lower}|>|\textit{upper}|$, so median $=5$.
  \item Add $15$: push to \textit{lower}$=\{5,15\}$ (max-heap top $15$).
        Pop max ($15$), push to \textit{upper}$=\{15\}$.
        $|\textit{upper}|=1=|\textit{lower}|=1$, no rebalance needed.
        Median: equal sizes, $(5+15)/2=10$.
  \item Add $1$: push to \textit{lower}$=\{1,5\}$ (top $5$). Pop max
        ($5$), push to \textit{upper}$=\{5,15\}$ (top $5$).
        $|\textit{upper}|=2>|\textit{lower}|=1$: pop \textit{upper}'s min
        ($5$), push to \textit{lower}$=\{1,5\}$ (top $5$).
        \textit{upper}$=\{15\}$. Median: $|\textit{lower}|=2>
        |\textit{upper}|=1$, median $=5$.
  \item Add $3$: push to \textit{lower}$=\{1,3,5\}$ (top $5$). Pop max
        ($5$), push to \textit{upper}$=\{5,15\}$ (top $5$).
        $|\textit{upper}|=2=|\textit{lower}|=2$, no rebalance. Median:
        equal sizes, $(\text{top of lower}=3 + \text{top of upper}=5)/2
        =4$.
\end{itemize}

Checking by hand: after $[5,15,1,3]$, sorted is $[1,3,5,15]$, median
$=(3+5)/2=4$. Correct.

\subsection*{C++ implementation}

\begin{lstlisting}
#include <bits/stdc++.h>
using namespace std;

class MedianFinder {
    priority_queue<int> lower; // max-heap: lower half
    priority_queue<int, vector<int>, greater<int>> upper; // min-heap: upper half

public:
    void addNum(int x) {
        lower.push(x);
        upper.push(lower.top());
        lower.pop();

        if (upper.size() > lower.size()) {
            lower.push(upper.top());
            upper.pop();
        }
    }

    double findMedian() {
        if (lower.size() > upper.size()) {
            return lower.top();
        }
        return (lower.top() + upper.top()) / 2.0;
    }
};

int main() {
    MedianFinder mf;
    for (int x : {5, 15, 1, 3}) {
        mf.addNum(x);
        cout << "Median after adding " << x << ": "
             << mf.findMedian() << endl;
    }
    return 0;
}
\end{lstlisting}

\begin{complexitybox}
\textbf{Time Complexity.} Each $\textit{addNum}$ performs a constant
number of heap pushes and pops, each $O(\log n)$, giving
\[
  O(\log n) \text{ per insertion}, \qquad O(1) \text{ per median query}.
\]
\textbf{Space Complexity.} $O(n)$ to store all elements seen so far,
split between the two heaps.
\end{complexitybox}

\begin{takeawaybox}
The two-heap pattern -- a max-heap for ``everything below the median'' and
a min-heap for ``everything above,'' kept balanced in size -- is the
standard way to maintain order statistics (median, or more generally any
fixed percentile with a fixed split ratio) over a stream that only grows.
Whenever a running statistic needs comparison to \emph{both} a lower and
an upper neighborhood simultaneously, suspect this pattern before reaching
for a single heap.
\end{takeawaybox}
